\chapter{Architecture}
\label{ch:architecture}

The admin bot is a \emph{maubot} plugin (\cmd{ch.unibas.dmi.marvin.adminbot}) written in
Python. maubot is a bot framework that runs plugins inside one process and gives each
running \emph{instance} a Matrix client, a database and optionally a small web app. The
plugin uses all three.

\section{Components}

\begin{center}
\begin{tikzpicture}[
  box/.style={draw=accent,fill=white,rounded corners=3pt,font=\footnotesize\sffamily,align=center,
    minimum height=1cm,text width=2.2cm},
  wide/.style={box,text width=6.0cm},
  zone/.style={draw=gray!60,dashed,rounded corners=6pt,inner sep=7pt},
  arr/.style={-{Stealth[length=2mm]},accent,thick},
  both/.style={{Stealth[length=2mm]}-{Stealth[length=2mm]},accent,thick},
  lbl/.style={font=\scriptsize\sffamily,text=gray,align=center}]
  \node[box] (element) at (-0.3,1.3) {Element\\(professors, students)};
  \node[box] (browser) at (-0.3,-0.5) {Browser\\(signup page)};
  \node[box] (nginx) at (2.9,0.4) {host nginx\\TLS, routing};
  \node[wide] (synapse) at (9.35,1.9) {Synapse \quad \cmd{http://synapse:8008}};
  \node[box,fill=accentlight] (main) at (7.3,-0.2) {maubot \enquote{main}\\\textbf{admin bot plugin}};
  \node[box] (userbots) at (11.4,-0.2) {maubot \enquote{userbots}\\professor bots};
  \node[wide] (pg) at (9.35,-2.3) {PostgreSQL: synapse, maubot, userbots};
  \draw[arr] (element.east) -- (nginx.west);
  \draw[arr] (browser.east) -- (nginx.west);
  \draw[arr] (nginx.east) -- node[lbl,above,sloped]{\cmd{/_matrix}} (synapse.west);
  \draw[arr] (nginx.east) -- node[lbl,below,sloped]{signup route} (main.west);
  \draw[both] (main.north) -- node[lbl,right,xshift=1pt]{client-server\\+ admin API} (main.north |- synapse.south);
  \draw[both] (userbots.north) -- node[lbl,left,xshift=-1pt]{client-server\\API} (userbots.north |- synapse.south);
  \draw[arr] (main.east) -- node[lbl,above]{management} node[lbl,below]{API} (userbots.west);
  \draw[both] (main.south) -- (main.south |- pg.north);
  \draw[both] (userbots.south) -- (userbots.south |- pg.north);
  \draw[both] (synapse.east) -- ++(0.55,0) |- (pg.east);
  \coordinate (edge) at ($(synapse.east)+(0.6,0)$);
  \begin{scope}[on background layer]
    \node[zone,fit=(synapse)(main)(userbots)(pg)(edge),
      label={[lbl]above:Docker network}] {};
  \end{scope}
\end{tikzpicture}
\end{center}

\begin{description}[style=nextline,leftmargin=1.5em,font=\normalfont\bfseries]
  \item[Synapse] The Matrix homeserver for \server. Federation is off.
  \item[host nginx] Terminates TLS. All of \cmd{/_matrix} goes to Synapse, except the prefix
    \cmd{/_matrix/maubot/plugin/<instance>/}, which goes to maubot so the signup page is
    reachable from outside. The maubot management API and UI are not routed; admins reach
    them through an SSH tunnel (\cmd{https://localhost:8443/_matrix/maubot/}).
  \item[maubot \enquote{main}] (container \cmd{matrix-maubot}) runs the admin bot instance
    with the client \cmd{@adminbot}. It reaches Synapse directly on the Docker network.
  \item[maubot \enquote{userbots}] (container \cmd{matrix-userbots}) runs the professor bots.
    It is a second, separate maubot without published ports; the admin bot manages it
    through its management API (\autoref{sec:why-userbots}).
  \item[PostgreSQL] Separate databases for Synapse, the main maubot (incl.\ the admin bot's
    tables) and userbots.
\end{description}

\section{Inside the plugin}

\begin{center}
\small
\begin{tabularx}{\textwidth}{@{}lL@{}}
\toprule
\textbf{Module} & \textbf{Responsibility} \\
\midrule
\file{bot.py} & plugin class: builds the services, startup self-check, cleanup timer, registers handlers \\
\file{commands/base.py} & \cmd{CommandRouter}, the only Matrix event handler; \cmd{@command}; shared helpers (replies, files, plans) \\
\file{commands/*.py} & one class per command group (\cmd{course}, \cmd{invite}, \cmd{groups}, \dots) \\
\file{authz.py} & roles and the single permission check \cmd{require(role, course)} \\
\file{confirm.py}, \file{jobs.py} & plans with confirm codes; resumable, throttled job runner \\
\file{signups.py}, \file{delivery.py} & signup and reset tokens, reset codes, pending invites; CSV/inline delivery \\
\file{commands/password.py} & \cmd{!password} and \cmd{!password code}, the public commands for every user \\
\file{notices.py} & messages held until the user has a device that can decrypt them \\
\file{web/} & signup, reset and recover pages, rate limiter, per-account lock for reset codes \\
\file{synapse_admin.py} & wrappers for the Synapse admin API \\
\file{matrix_ops.py} & rooms, spaces, power levels, membership, DM check, sending \\
\file{templates.py} & template schema, parse/dump, snapshot of a live course, apply \\
\file{userbots.py} & \file{.mbp} validation, client for the userbots maubot, provisioning \\
\file{identity.py}, \file{parsing.py}, \file{roster.py} & pure logic: email to user ID, arguments and CSV, roster diff \\
\file{db.py}, \file{migrations.py} & SQL (SQLite in tests, PostgreSQL in production) \\
\file{audit.py} & audit table and audit room \\
\file{strings/en.py} & every user-facing text, including the detailed help \\
\bottomrule
\end{tabularx}
\end{center}

\section{The path of a command}
\label{sec:command-path}

\begin{enumerate}
  \item The professor's Element encrypts the message and sends it to Synapse.
  \item maubot's client for \cmd{@adminbot} receives it through its \cmd{/sync} loop, and
        mautrix decrypts it.
  \item \cmd{CommandRouter.on_message} checks, cheapest first:
        \begin{itemize}
          \item the message is text and starts with \cmd{!<name>};
          \item the sender is local and \emph{known} (admin, professor or staff of a course).
                Other local users (students) may only use the \emph{public} commands
                (\cmd{!help}, \cmd{!password}, \cmd{!password code}), and only in a DM;
                everything else they send, and anything outside a DM, is ignored without
                a reply;
          \item the room is a DM between exactly the bot and the sender
                (\autoref{sec:dm-rule}). Otherwise known users get \enquote{please DM me},
                and unknown command names are ignored, because other bots in course rooms
                use \cmd{!} commands too;
          \item the startup checks the command depends on have passed.
        \end{itemize}
  \item The router reacts with \robot, records the DM for later notifications, and calls the
        handler. \cmd{!x help}, \cmd{!x y --help} and unknown subcommands are answered by the
        help handler.
  \item The handler parses the arguments, calls \cmd{authz.require(role, course)}, and does
        the work through the services.
  \item Errors become replies: \cmd{AdminBotError} subclasses carry a strings key, Matrix
        errors are shown with the server's reason, anything else as an internal error
        (with the traceback in the log).
  \item The router logs how long the message took to reach the bot and how long the handler
        took (\autoref{sec:timing}).
\end{enumerate}

\section{Plans and jobs}

Bulk or destructive commands never act directly. The handler builds a \emph{plan}: a
summary, a list of items (one per user or room) and a context (course, the role to
re-check, flags). \cmd{ConfirmRegistry.propose} stores it in the \cmd{plan} table with a
random four-character code and a ten-minute expiry; a user has at most one open plan.

\cmd{!confirm <code>} pops the plan (only for its requester), re-checks the role against
the course and hands it to the \cmd{JobRunner}:

\begin{itemize}
  \item the job and one \cmd{job_item} row per item are written to the database;
  \item an optional \cmd{on_start} hook runs once (e.g.\ \cmd{!invite} creates the signup
        links and delivers the CSV);
  \item a single worker task processes items one by one, with \cmd{jobs.delay_seconds}
        between them, and edits one progress message every few seconds;
  \item item failures are recorded per item and don't stop the job;
  \item at the end, an optional \cmd{on_finish} hook runs, then the report (with failed
        items and a next-step hint) and an audit entry.
\end{itemize}

Unfinished jobs are re-queued when the plugin starts, so a restart resumes where it stopped.

\section{Files}

Commands that need a file either use the file the command replies to, or register a
\emph{pending upload} (in memory, per room and sender, \cmd{uploads.wait_minutes}) and take
the next file the sender posts in that DM. Encrypted attachments are decrypted with
\cmd{mautrix.crypto.attachments}. If an upload handler raises a \cmd{ValidationError}, the
router registers the same handler and context again, so the user only re-sends the file.
